\documentclass[11pt,a4paper]{report}
\usepackage[utf8]{inputenc}
\usepackage[T1]{fontenc}
\usepackage[spanish]{babel}
\usepackage[margin=2.5cm]{geometry}
\usepackage{lmodern}
\usepackage{microtype}
\usepackage{tabularx}
\usepackage{hyperref}
\hypersetup{hidelinks,pdftitle={ServerFort: Memoria inicial},pdfauthor={Hugo Poblaciones, Sergi Blanca y Unai Hervella}}
\setlength{\parindent}{0pt}
\setlength{\parskip}{0.6em}
\begin{document}
\begin{titlepage}
\centering
{\large Proyecto de Grado Superior ASIX\par}
\vspace{1cm}
{\large Memoria inicial (anteproyecto)\par}
\vspace{2cm}
{\Huge\bfseries ServerFort\par}
\vspace{1cm}
{\Large Infraestructura Ubuntu Server segura, monitorizada y automatizada\par}
\vspace{0.8cm}
{\large De servidor Linux a infraestructura empresarial completa\par}
\vfill
Hugo Poblaciones\\Sergi Blanca\\Unai Hervella
\vspace{1cm}

Proyecto Optativo -- ASIX\\Profesor: Jass Riahi
\vspace{1cm}

19 de octubre de 2026
\end{titlepage}
\tableofcontents

\chapter{Introducción}

\section{De qué va el proyecto}
 ServerFort consiste en montar una infraestructura de empresa pequeña con Ubuntu Server: servidores, red, seguridad y copias de seguridad. La idea es que, una vez montada, se mantenga casi sola y que nadie tenga que estar todo el día pendiente. 
\section{El problema}
 Muchas pymes ponen un servidor Linux sin un firewall bien configurado, sin backups automáticos y sin nadie que vigile si algo falla. Como normalmente no tienen un administrador de sistemas a jornada completa, los problemas se detectan tarde. Si entra alguien que no debe o se estropea un disco, ya se ha perdido tiempo o datos. Un fallo así puede parar el negocio varios días, por eso creemos que merece la pena resolverlo. 
\section{Para quién está pensado}
 
\par\noindent\textbullet\quad  Una pyme o un departamento de IT pequeño que no tiene un administrador de sistemas a tiempo completo. 
\par\noindent\textbullet\quad  Una organización que necesite red, servicios (web, archivos, DNS...) y acceso remoto seguro sin depender de terceros. 
\par\noindent\textbullet\quad  Nosotros mismos, ya que nos sirve para poner en práctica lo aprendido en el ciclo. 
\section{Objetivos}
 Queremos pasar de un Ubuntu Server sin configurar a una pequeña infraestructura de empresa. Para ello vamos a: 
\par\noindent\textbullet\quad  Diseñar una red separada por zonas, con DMZ, servidores y acceso remoto. 
\par\noindent\textbullet\quad  Montar los servicios: DNS, DHCP, Apache, archivos y VPN. 
\par\noindent\textbullet\quad  Proteger el sistema con firewall, permisos y auditoría (y IDS/IPS si nos da tiempo). 
\par\noindent\textbullet\quad  Controlar el estado de las máquinas con Prometheus, logs y alertas. 
\par\noindent\textbullet\quad  Automatizar tareas con Ansible y scripts. 
\par\noindent\textbullet\quad  Hacer copias de seguridad y probar a restaurarlas, con un plan de recuperación ante desastres. 
\section{Qué aporta de diferente}
 No nos queremos quedar solo en instalar un servidor. Vamos a simular ataques y fallos (un escaneo con Nmap, un servicio caído, un disco lleno) para comprobar que la seguridad y las alertas funcionan de verdad. Además, todo se hará con Ansible para poder repetir la instalación desde cero en pocos minutos, y juntamos redes, seguridad, monitorización y recuperación en el mismo proyecto. 
\section{Estado del proyecto}
 Esta memoria se ha escrito antes de empezar el proyecto, así que todavía no hay resultados. Aquí explicamos qué queremos hacer, cómo lo vamos a organizar y cómo comprobaremos que funciona. Todo lo que aparece es una previsión y puede cambiar durante el curso. Las decisiones concretas (programas que usaremos, direcciones IP, configuraciones) las iremos tomando en cada sprint. Los resultados, las pruebas y las conclusiones se añadirán en la memoria final.

\chapter{Alcance}
\begin{tabularx}{\textwidth}{|p{3.2cm}|X|}
\hline
\textbf{Parte} & \textbf{Qué incluye} \\ \hline
Imprescindible (MVP) & Red separada, DNS/DHCP, servidor web, firewall, backups y monitorización básica con Prometheus y Grafana. \\ \hline
Ampliaciones & Suricata (IDS/IPS), LDAP/FreeIPA y servicios en contenedores Docker. \\ \hline
Fuera del alcance & Programar aplicaciones propias y montarlo en hardware real. Todo se hace con máquinas virtuales. \\ \hline
\end{tabularx}
\section{Máquinas virtuales}
 Usaremos dos máquinas virtuales: 
\par\noindent\textbullet\quad  Ubuntu Server, donde están los servicios. 
\par\noindent\textbullet\quad  Ubuntu Client, desde donde hacemos todas las pruebas contra el servidor y nos conectamos a él de forma remota.

\chapter{Arquitectura}
La infraestructura se divide en estas partes: 
\par\noindent\textbullet\quad  Internet : entrada y salida de tráfico. 
\par\noindent\textbullet\quad  Firewall : NAT, reglas, VPN y separación de redes. 
\par\noindent\textbullet\quad  DMZ : web y servicios que se publican. 
\par\noindent\textbullet\quad  LAN / usuarios : acceso a los recursos internos. 
\par\noindent\textbullet\quad  Servidores : DNS, LDAP, archivos, monitorización y backup. 
\par\noindent\textbullet\quad  Laboratorio : pruebas de seguridad controladas. El esquema de red se hará en el Sprint 1 y se añadirá a la memoria final.

\chapter{Red (previsión)}
Vamos a dividir la red en cinco VLAN para separar cada función y que, si pasa algo en una zona, no afecte a todas. 

\begin{tabularx}{\textwidth}{|r|l|X|l|}
\hline
\textbf{VLAN} & \textbf{Nombre} & \textbf{Para qué sirve} & \textbf{Red} \\ \hline
10 & Administración & Gestionar equipos y servidores & Por definir \\ \hline
20 & Usuarios & Acceso a recursos de la empresa & Por definir \\ \hline
30 & Servidores & Servicios internos & Por definir \\ \hline
40 & DMZ & Servicios expuestos & Por definir \\ \hline
50 & Laboratorio & Pruebas y pentesting & Por definir \\ \hline
\end{tabularx}

Usaremos IPv4, DHCP, DNS y reglas de firewall entre las zonas. La idea es bloquear todo por defecto y abrir solo lo necesario. Las subredes y las reglas concretas entre VLAN se definirán en el Sprint 1 y el Sprint 3.

\chapter{Servicios}

\par\noindent\textbullet\quad  DNS y DHCP : resolución de nombres y configuración automática de la red. Los programas concretos se elegirán en el Sprint 2. 
\par\noindent\textbullet\quad  SSH : administración remota con claves. 
\par\noindent\textbullet\quad  Web : Apache con HTTPS, en la DMZ. El dominio y el certificado se decidirán en el Sprint 2. 
\par\noindent\textbullet\quad  Archivos : Samba con permisos por usuario y grupo, para que cada usuario solo vea sus cosas. 
\par\noindent\textbullet\quad  VPN : WireGuard para conectarse desde fuera de forma segura. 
\section{Configuración}
 La configuración de cada servicio se documentará durante los Sprints 2 y 3 y se incluirá en la memoria final.

\chapter{Seguridad}
Medidas que vamos a aplicar: 
\par\noindent\textbullet\quad  SSH solo con clave y sin acceso directo como root. 
\par\noindent\textbullet\quad  Instalar solo los servicios necesarios y mantener el sistema actualizado. 
\par\noindent\textbullet\quad  Usar sudo, grupos y permisos, con política de contraseñas. 
\par\noindent\textbullet\quad  Firewall con puertos mínimos y reglas por zonas. 
\par\noindent\textbullet\quad  Suricata para detectar tráfico sospechoso (ampliación). La herramienta de firewall y su configuración se decidirán en el Sprint 3.

\chapter{Monitorización}

\par\noindent\textbullet\quad  Prometheus recoge métricas de los servidores y servicios. 
\par\noindent\textbullet\quad  Grafana muestra gráficas de CPU, RAM, disco, red y disponibilidad. 
\par\noindent\textbullet\quad  Alertmanager avisa cuando hay una incidencia. 
\par\noindent\textbullet\quad  Los logs se centralizan para poder revisarlos. Los paneles y las alertas concretas se definirán en el Sprint 4.

\chapter{Copias de seguridad}
Seguiremos la regla 3-2-1: tres copias, en dos tipos de soporte y una de ellas fuera. Las copias serán diarias, semanales y mensuales, con un tiempo de retención. Se comprimen, se cifran y se guardan en un sitio distinto al servidor. También se comprobará automáticamente que las copias están bien, y probaremos a borrar datos del laboratorio y recuperarlos. La herramienta, lo que se copia y el tiempo de retención se definirán en el Sprint 5. 
\section{Plan de recuperación (Disaster Recovery)}
 El plan de recuperación se redactará en el Sprint 5. Este plan consistirá en un plan por si una de las máquinas virtuales se rompe, recuperarla a través de una copia.

\chapter{Automatización}

\par\noindent\textbullet\quad  Ansible : usuarios, paquetes, SSH, firewall y servicios. 
\par\noindent\textbullet\quad  Scripts : backups, comprobaciones y mantenimiento. 
\par\noindent\textbullet\quad  Docker/Podman : para algunos servicios (ampliación). Para montar todo en máquinas nuevas, se ejecutaría el playbook principal: \par\texttt{ansible-playbook -i inventory/hosts.ini site.yml}\par La estructura del repositorio se definirá al empezar el proyecto.

\chapter{Pruebas de seguridad}
Haremos las pruebas siguiendo este ciclo: ataque o fallo  $\rightarrow$  detección  $\rightarrow$  corrección  $\rightarrow$  validación (comprobar otra vez que está arreglado). 
\par\noindent\textbullet\quad  Servidor caído. 
\par\noindent\textbullet\quad  Disco lleno. 
\par\noindent\textbullet\quad  Servicio parado. 
\par\noindent\textbullet\quad  Intentos de acceso por SSH. 
\par\noindent\textbullet\quad  Escaneo de puertos con Nmap. Los resultados de las pruebas se añadirán en la memoria final.

\chapter{Organización}

\section{Reparto de tareas}
 
\par\noindent\textbullet\quad  Hugo Poblaciones : redes, firewall y VLAN. 
\par\noindent\textbullet\quad  Sergi Blanca : servicios (DNS, web, archivos) y Ansible. 
\par\noindent\textbullet\quad  Unai Hervella : seguridad, monitorización y backups. Los roles pueden cambiar durante el curso según lo que necesitemos. 
\section{Sprints}
\begin{enumerate}

\item Diseño de la red, las VLAN y la arquitectura.
\item Servicios base: DNS, DHCP, web y archivos.
\item Firewall, VPN y hardening.
\item Monitorización, dashboards y alertas.
\item Backups y pruebas de ataques simulados.
\item Documentación, pruebas finales y presentación.
\end{enumerate}
\chapter{Problemas previstos}

\par\noindent\textbullet\quad  Problemas técnicos : Puede que algún servicio no arranque bien en la VM y nos lleve tiempo arreglarlo. Dejamos margen en cada sprint. 
\par\noindent\textbullet\quad  Cosas nuevas : Ansible y Suricata no los hemos usado antes. Empezaremos con tutoriales antes de aplicarlos. 
\par\noindent\textbullet\quad  Falta de tiempo : Si nos retrasamos, lo primero que quitaremos será el laboratorio de pentesting, para no afectar a lo imprescindible. Los problemas reales que aparezcan se recogerán en la memoria final.

\chapter{Entrega y comprobación}

\section{Qué entregamos}
 Las máquinas virtuales con la red y los servicios, el repositorio (Ansible, scripts, configuraciones y documentación), las pruebas de seguridad, los dashboards y alertas, los backups con su restauración, y esta memoria con la demo. 
\section{Cómo demostramos que funciona}
\begin{enumerate}

\item Un usuario se autentica y solo accede a sus recursos.
\item Los servicios responden y están monitorizados.
\item El firewall bloquea el tráfico no autorizado.
\item Los intentos sospechosos generan registros y alertas.
\item Se hace un backup y se restaura.
\item Una incidencia simulada se detecta y se documenta.
\item La configuración se puede repetir con la automatización.
\end{enumerate}
Las capturas y evidencias de cada punto se recogerán durante el proyecto y se incluirán en la memoria final.\chapter{Resultados esperados y conclusiones}

\section{Qué esperamos conseguir}
 Al terminar el proyecto, esperamos tener una infraestructura funcionando que sea segura, esté vigilada, se pueda reconstruir con Ansible y permita recuperar los datos si algo falla. También esperamos aprender a defender, vigilar y recuperar un servidor, que es lo que se le pide a un administrador de sistemas en una empresa real. 
\section{Conclusiones}
 Las conclusiones se escribirán al terminar el proyecto, cuando tengamos resultados reales.

\end{document}
